\documentclass[11pt]{article}
\usepackage[margin=1in]{geometry}
\usepackage{amsmath,amssymb,amsthm}
\usepackage{graphicx}
\usepackage{hyperref}
\usepackage{xcolor}
\usepackage{enumitem}
\usepackage{booktabs}

\newtheorem{theorem}{Theorem}
\newtheorem{lemma}{Lemma}
\newcommand{\Lp}{L^p[0,1]}
\newcommand{\Dcal}{\mathcal D}
\newcommand{\Ecal}{\mathbb E}
\newcommand{\one}{\mathbf 1}

\title{A Haar-Rearrangement Solution of the Multiplicity-Two Problem\\
for Multiplication on $L^p[0,1]$}
\author{Candidate solution packet for arXiv:1707.04798, Problem 1}
\date{August 9, 2026}

\begin{document}
\maketitle

\begin{center}
\fcolorbox{black}{gray!10}{\parbox{0.91\textwidth}{
\textbf{Status.} Candidate full solution; verdict \textbf{likely valid}, pending
independent human review. The proof gives an affirmative answer for every
$1<p<\infty$, including $p=2$.
}}
\end{center}

\section{Source question}

Let $M$ be multiplication by the coordinate on $\Lp$,
\[
  (Mf)(t)=t f(t).
\]
Problem 1 in Section 4 of Boedihardjo's paper \cite{Boedihardjo} asks, for
$1<p<\infty$ and $p\ne2$, whether there is a compact operator $K$ on $\Lp$
such that $M\oplus M$ is similar to $M+K$.

\begin{figure}[ht]
  \centering
  \includegraphics[width=\textwidth]{figures/open_problem_crop.png}
  \caption{Problem 1 on PDF page 9 of the source paper.}
\end{figure}

The source notes on the following page that the answer is known to be
affirmative when $p=2$.

\section{Proof intuition}

The coordinate multiplier is compactly close to a diagonal operator on the
standard Haar basis: attach to the Haar function on a dyadic interval $I$ the
center $c_I$ of that interval. To fit two copies of this diagonal into one
copy, send the Haar coordinate indexed by $I$ in the first copy to the
coordinate indexed by the left child $I0$, and send the corresponding
coordinate in the second copy to the right child $I1$.

There are two points to prove. First, each fixed-child rearrangement is an
isomorphism onto its Haar-coordinate range. The upper estimate follows
pointwise from Haar square functions; the lower estimate follows from Stein's
vector-valued conditional-expectation inequality. Second, after the
rearrangement, the value $c_I$ is placed on $I0$ and $I1$, whereas the original
diagonal places $c_{I0}$ and $c_{I1}$ there. The errors have magnitude
$|I|/4$, hence tend uniformly to zero by Haar level and define a compact
diagonal operator.

\section{Dyadic preliminaries}

Let $\Dcal$ be the dyadic intervals in $[0,1]$. For $I\in\Dcal$, let $I0$
and $I1$ denote its left and right children, let $c_I$ be its center, and use
the $L^p$-normalized Haar function
\[
  h_I=|I|^{-1/p}(\one_{I0}-\one_{I1}).
\]
Together with $e=\one_{[0,1]}$, the family
$\{e\}\cup\{h_I:I\in\Dcal\}$ is an unconditional Schauder basis of $\Lp$.

For a finite Haar expansion $x=\sum_I a_Ih_I$, set
\[
  Sx(t)=\left(\sum_I |a_I|^2 |I|^{-2/p}\one_I(t)\right)^{1/2}.
\]
The martingale square-function theorem \cite{Stein} gives constants depending only on $p$
such that
\begin{equation}\label{eq:square}
  A_p\|x\|_p\le \|Sx\|_p\le B_p\|x\|_p.
\end{equation}
We also use Stein's vector-valued inequality for dyadic conditional
expectations \cite{Stein}: if $\Ecal_n$ is conditional expectation onto the
dyadic sigma-algebra at level $n$, then
\begin{equation}\label{eq:stein}
 \left\|\left(\sum_n|\Ecal_n f_n|^2\right)^{1/2}\right\|_p
 \le C_p
 \left\|\left(\sum_n|f_n|^2\right)^{1/2}\right\|_p.
\end{equation}

\section{The child rearrangements}

Let $H$ be the closed span of the Haar functions $h_I$.

\begin{lemma}\label{lem:child}
For $\varepsilon\in\{0,1\}$, the rule
\[
  T_\varepsilon h_I=h_{I\varepsilon}
\]
extends to an isomorphism from $H$ onto the closed Haar-coordinate subspace
\[
  Y_\varepsilon=\overline{\operatorname{span}}
  \{h_{I\varepsilon}:I\in\Dcal\}.
\]
\end{lemma}

\begin{proof}
We prove the claim first on finite Haar sums. If
$x=\sum_I a_Ih_I$, then
\[
 S(T_\varepsilon x)(t)^2
 =2^{2/p}\sum_I|a_I|^2|I|^{-2/p}\one_{I\varepsilon}(t)
 \le 2^{2/p}Sx(t)^2.
\]
Together with \eqref{eq:square}, this proves the upper bound for
$T_\varepsilon$.

For the reverse bound, for every level $n$ define
\[
 f_n=\sum_{\substack{I\in\Dcal\\ |I|=2^{-n}}}
 a_I|I|^{-1/p}\one_{I\varepsilon}.
\]
The intervals at a fixed level are disjoint, and
\[
 \Ecal_n f_n=\frac12
 \sum_{|I|=2^{-n}}a_I|I|^{-1/p}\one_I.
\]
Consequently,
\[
 Sx=2\left(\sum_n|\Ecal_n f_n|^2\right)^{1/2},
 \qquad
 \left(\sum_n|f_n|^2\right)^{1/2}
 =2^{-1/p}S(T_\varepsilon x).
\]
Applying \eqref{eq:stein}, followed by \eqref{eq:square} twice, gives
\[
 \|x\|_p\le C'_p\|T_\varepsilon x\|_p.
\]
Thus $T_\varepsilon$ is bounded and bounded below. Its range is therefore
closed and is exactly $Y_\varepsilon$.
\end{proof}

Define $U:(\Lp\oplus\Lp)_p\to\Lp$ by
\begin{equation}\label{eq:U}
 U(ae+x,be+y)=ae+b h_{[0,1]}+T_0x+T_1y,
 \qquad x,y\in H.
\end{equation}
The four coordinate families in \eqref{eq:U} are disjoint, and together they
exhaust the target Haar basis: the constant $e$, the root Haar function
$h_{[0,1]}$, all left-child Haar functions, and all right-child Haar
functions. Since the Haar basis is unconditional, the coordinate projections
onto these four families are bounded. Lemma \ref{lem:child} therefore shows
that $U$ is bounded, onto, and has a bounded inverse.

\section{A compact Haar diagonalization}

Define a diagonal operator $D$ on the Haar basis by
\begin{equation}\label{eq:D}
  De=\frac12e,
  \qquad
  Dh_I=c_Ih_I \quad(I\in\Dcal).
\end{equation}

\begin{lemma}\label{lem:compactdiag}
$C=M-D$ is compact on $\Lp$.
\end{lemma}

\begin{proof}
Let $P_n$ be the Haar-coordinate projection onto the functions $h_I$ with
$|I|=2^{-n}$. These projections are uniformly bounded by unconditionality.
If $x=\sum_{|I|=2^{-n}}a_Ih_I$, the supports of the summands are disjoint and
$|t-c_I|\le |I|/2$ on $I$. Hence
\[
 \|(M-D)x\|_p\le 2^{-n-1}\|x\|_p.
\]
Let $Q_N$ be the coordinate projection onto the constant and the Haar levels
at most $N$. It follows that
\[
 \left\|(M-D)(I-Q_N)\right\|
 \le C_p\sum_{n>N}2^{-n-1}\longrightarrow0.
\]
Since $(M-D)Q_N$ has finite rank, $M-D$ is the operator-norm limit of
finite-rank operators and is compact.
\end{proof}

\section{Full solution}

\begin{theorem}\label{thm:main}
Let $1<p<\infty$ and let $M$ be multiplication by $t$ on $L^p[0,1]$. There
are an isomorphism
\[
 U:(L^p[0,1]\oplus L^p[0,1])_p\longrightarrow L^p[0,1]
\]
and a compact operator $K$ on $L^p[0,1]$ such that
\[
  U(M\oplus M)U^{-1}=M+K.
\]
In particular, Problem 1 of \cite{Boedihardjo} has an affirmative answer.
\end{theorem}

\begin{proof}
Use the isomorphism $U$ from \eqref{eq:U} and the diagonal $D$ from
\eqref{eq:D}. Put
\[
  A=U(D\oplus D)U^{-1}.
\]
This is diagonal in the target Haar basis. The two constant coordinates give
\[
  Ae=\tfrac12e,
  \qquad Ah_{[0,1]}=\tfrac12h_{[0,1]},
\]
and for every dyadic $I$ and $\varepsilon\in\{0,1\}$,
\[
  Ah_{I\varepsilon}=c_Ih_{I\varepsilon}.
\]
On the other hand, $Dh_{I\varepsilon}=c_{I\varepsilon}h_{I\varepsilon}$, and
\[
 |c_I-c_{I\varepsilon}|=\frac{|I|}{4}.
\]
Thus $E=A-D$ is diagonal on an unconditional basis and its diagonal
coefficients tend to zero. Its truncations converge in operator norm, so $E$
is compact.

By Lemma \ref{lem:compactdiag}, $M=D+C$ for a compact $C$. Therefore
\begin{align*}
 U(M\oplus M)U^{-1}
 &=U(D\oplus D)U^{-1}+U(C\oplus C)U^{-1}\\
 &=D+E+U(C\oplus C)U^{-1}\\
 &=M+\bigl(E+U(C\oplus C)U^{-1}-C\bigr).
\end{align*}
The operator in parentheses is compact. Taking it to be $K$ proves the
claim.
\end{proof}

\section{Verification notes}

The main possible failure modes were checked as follows.
\begin{enumerate}[leftmargin=*]
\item The proof does not infer compactness merely from small errors on basis
vectors. Lemma \ref{lem:compactdiag} gives a summable operator-norm tail, and
$E$ is a null diagonal on an unconditional basis.
\item The map $I\mapsto I\varepsilon$ does alter nesting, so arbitrary-basis
permutability cannot be invoked. Lemma \ref{lem:child} proves the two required
rearrangements directly.
\item The two image subspaces need not be an isometric $\ell^p$ sum. Bounded
coordinate projections give the needed topological direct sum, and finite
product norms are equivalent.
\item The finite coordinates match: the first constant maps to the target
constant, the second constant to the root Haar function, and all remaining
target Haar functions are uniquely a left or right child.
\end{enumerate}

The script \texttt{code/verify\_dyadic\_absorption.py} checks finite dyadic
truncations, exact coordinate coverage, parent-child center errors, and random
$L^p$ norm ratios. These computations are sanity checks and are not used as
proof. With 500 fixed-seed samples at every depth through 8, the exact checks
passed and the randomized ratios, over $p\in\{1.2,1.5,3,5\}$, lay between
$0.644458$ and $1.687973$.

\section{Novelty and limitations}

The lightweight run indexes were searched using arXiv:1707.04798, the exact
problem language, multiplicity-two terminology, compact perturbation, and
Haar-rearrangement phrases. A bounded external search used the exact question,
the paper title and author, the arXiv id, and close variants. It found the
source arXiv and journal records but no later work claiming this resolution.
The search was bounded rather than exhaustive, so the mathematical novelty
claim has moderate confidence.

The theorem resolves Problem 1 and supplies an explicit Haar-coordinate
similarity. It does not address Problem 2 in the source paper, which asks for a
1-summing perturbation for a multiplication operator on a square in the
complex plane.

\section{Human-review recommendation}

The packet should be reviewed as a candidate full solution. The highest-value
checks are the normalization and applicability of Stein's inequality in
Lemma \ref{lem:child}, followed by the finite-coordinate bookkeeping in
\eqref{eq:U}. The compactness and conjugation steps are then routine.

\begin{thebibliography}{9}

\bibitem{Boedihardjo}
M.~T. Boedihardjo,
\emph{Multiplication operators on $L^p$},
Studia Math. 244 (2019), 309--319;
\href{https://arxiv.org/abs/1707.04798}{arXiv:1707.04798}.

\bibitem{Stein}
E.~M. Stein,
\emph{Topics in Harmonic Analysis Related to the Littlewood--Paley Theory},
Annals of Mathematics Studies 63, Princeton University Press, 1970.

\end{thebibliography}

\end{document}
